\documentclass{article}
\usepackage{styles/project_style} % Assuming you have this style file



% Page Geometry
\geometry{a4paper, margin=1in}

% Title, Author, and Date (customize as needed)
\title{Course Grading Guidelines}
\author{Václav Petrák} % Placeholder
\date{\today}

\begin{document}
\maketitle

\section*{Classification and Grading}

There is no exam. Your grade will be based on your independent work. There are three outputs used for grading.

\begin{enumerate}
    \item     \textbf{Journal Club Presentation}:  30\% of the final grade
    \item \textbf{Capstone Project Presentation}:
    { 40\% of the final grade}
    \item  \textbf{Code Review Discussion}:
    { 30\% of the final grade}
\end{enumerate}
 
   



\section{Journal Club}
A journal club is a group who meet to critically evaluate articles in the academic literature. Journal clubs:
\begin{itemize}
    \item Encourages critical thinking.
    \item Helps in understanding the studied concepts and methods.
    \item Helps keep up with the latest research.
    \item Improves communication and presentation skills.
    \item Presentation simulates experience of a scientist.
\end{itemize}
\textbf{The main purpose in this course is to explore how methods we learned are applied in current research.}

\subsection{Your Task for Journal Club}
\begin{itemize}
    \item You will read a research paper related to the topic of your capstone project.
    \item Your task is to read the paper and prepare a 15-minute presentation, followed by a 5-minute discussion, to explain the paper.
    \item The presentation dates will be arranged individually.
\end{itemize}

\subsection{Journal Club Presentation Content}
The Journal Club presentation will cover:
{ % Group to apply  to the list

\begin{itemize}
    \item \textbf{Introduction:} Provide a High-level Overview on the background of the research topic and the main objectives of the paper.
    \item \textbf{Methods:} Give overview of the methodology used in the paper.
    \item \textbf{Results:} Highlight the key results and conclusions of the paper.
    \item \textbf{Discussion:} Discuss the practical implications of the research and its applications in real-world situations. Discuss any limitations of the paper.
    \item \textbf{Conclusions:} Summarize the main contributions and overall findings of \textbf{the paper}. Reiterate its key takeaways and their broader implications.
    \item \textbf{Q\&A:} Answer questions from the audience (lecturer).
\end{itemize}
}

\subsection{Grading Criteria for Journal Club Presentation}
\textbf{Presentation Design and Delivery (50\%)}
{ % Group to apply  to the list
\begin{itemize}
    \item \textbf{Design:} The presentation is well-organized and easy to follow.
    \item \textbf{Timing:} The presenter adhered to the 15-minute presentation format.
    \item \textbf{Delivery:} Clear and confident speech with appropriate pacing.
    \item \textbf{Questions:} Convincing response to questions.
\end{itemize}
}



\textbf{Content of the presentation (50\%)}
{ % Group to apply  to the list
\begin{itemize}
    \item \textbf{Introduction:} The introduction provides a high-level overview of the topic: explains the research motivation and the research question.
    \item \textbf{Methods:} Explains the experimental approach used by the authors. Unfamiliar methods are explained without going into excessive detail.
    \item \textbf{Results:} The key results from the paper are explained.
    \item \textbf{Critical evaluation:} For example critique of methods, ideas for improvement, significance of the results.
\end{itemize}
}

\section{Capstone Project}
The Capstone project is your independent project.
\begin{itemize}
    \item You will work on an independent project that builds on and further extends course content.
    \item Output of your work will be a presentation and a code.
    \item The presentation dates will be arranged individually.
\end{itemize}


\subsection{Capstone Project Presentation Content}
The Capstone Project presentation will cover:
{ % Group to apply  to the list
\begin{itemize}
    \item \textbf{Introduction:} Provide a High-level Overview on the background of the research topic and the main objectives of \textbf{your project}.
    \item \textbf{Methods:} Give overview of the methodology used in \textbf{your project}.
    \item \textbf{Results:} Highlight the key results and conclusions of \textbf{your project}.
    \item \textbf{Discussion:} Discuss the practical implications of \textbf{your project's findings} and their applications in real-world situations. Discuss any limitations of \textbf{your project}.
    \item \textbf{Conclusions:} Summarize the main contributions and overall findings of \textbf{your project}. Reiterate its key takeaways and their broader implications.
    \item \textbf{Q\&A:} Answer questions from the audience (lecturer).
\end{itemize}
}

\subsection{Grading Criteria for Capstone Project}
\textbf{Presentation Design and Delivery (50\%)}
{ % Group to apply  to the list
\begin{itemize}
    \item \textbf{Design:} The presentation is well-organized and easy to follow.
    \item \textbf{Plots:} Clear graphical presentation of simulation results that address project objectives and show model behavior.
    \item \textbf{Timing:} The presenter adhered to the 15-minute presentation format.
    \item \textbf{Delivery:} Clear and confident speech with appropriate pacing.
    \item \textbf{Questions:} Convincing response to questions.
\end{itemize}
}


\textbf{Model Implementation (50\%)}
{ % Group to apply  to the list
\begin{itemize}
    \item \textbf{Correctness:} Correct function of fundamental models.
    \item \textbf{Understanding:} Demonstrated understanding of core concepts.
    \item \textbf{Extension beyond core:} Successful and meaningful implementation of additional features (if the task specifies them).
    \item \textbf{Discussion:} Insightful discussion on how parameters/initial conditions affect outcomes. Identification and explanation of behaviors.
\end{itemize}
}

\section{Code Review}
Code review is a software quality assurance process where individuals examine source code to identify errors, enhance quality and maintainability, ensure adherence to standards, and promote knowledge sharing. Its aim is to improve the code before it's finalized. For students, code review is an effective learning tool. It helps to:
\begin{itemize}
    \item \textbf{Deepen conceptual understanding:} Analyzing and discussing code improves understanding of programming concepts and helps identify learning gaps.
    \item \textbf{Develop professional competencies:} It is building skills in receiving feedback, and articulating technical decisions.
    \item \textbf{Helps further learning:} It creates an environment for sharing knowledge, learning different approaches, and gaining insights from peers.
\end{itemize}
In this course, the code review specifically assesses your understanding of the capstone project, the quality of your code, and your proficiency in discussing its design and implementation.

\subsection{Code Review Content}
During Code Review we will discuss your code. Code Review assesses \textbf{your understanding}, quality, design, and documentation of the Capstone Project code.

{ % Group to apply  to the list
\begin{itemize}
    \item \textbf{Primary Focus:} The review concentrates on code for the Capstone Project.
    \item \textbf{Interactive Session:} Review involves an interactive code walk-through you will explain your code's functionality and design choices.
    \item \textbf{Demonstrating Ownership:} This is an opportunity for you to demonstrate your good understanding of the project, justify the decisions made during development, and articulate the purpose of different code segments.
    \item \textbf{Goal:} The ultimate goal is to ensure the correctness of your project, while demonstrating your grasp of the implemented concepts.
\end{itemize}
}

\subsection{Grading Criteria for Code Review}
\textbf{Readability and Maintainability (20\%)}
{ % Group to apply  to the list

\begin{itemize}
    \item Clean, well-formatted, and easy to follow code.
    \item Consistent naming conventions and style.
    \item Logical organization: functions that brake code into manageable components in neceessary.
    \item Sufficient and meaningful comments for complex/non-obvious code.
\end{itemize}
}


\textbf{Code understanding (80\%)}
{ % Group to apply  to the list

\begin{itemize}
    \item Ability to explain how any part of the code works.
    \item Ability to clearly articulate design decisions and justify coding choices made during development.
    \item Ability to suggest improvements, extensions or alternatives.
\end{itemize}
}

\end{document}